% ===========================================================================
% §Validation, limites et avertissement
% ===========================================================================
\section{Validation : la parité au cent près}
\label{sec:validation}

\subsection{Méthode}

La règle de travail du projet : \emph{ne rien présumer}. Chaque valeur
provient d'une source officielle citée ; chaque règle est tracée dans le
calculateur officiel du MFQ ; et l'ensemble est validé par une suite de
tests automatisés qui compare les sorties du moteur à celles de la
référence.

Concrètement, la suite de tests charge le code source du calculateur
officiel (tel que distribué par le site du MFQ), l'exécute dans un bac à
sable isolé, et confronte chaque sortie --- le revenu disponible, mais
aussi chacun des postes --- aux résultats du moteur TypeScript, \emph{au
cent près} :

\begin{itemize}
  \item un \textbf{balayage dense} : pour chacune des cinq situations, le
        revenu varie de 0 à \D{150000} par pas de \D{2000}, avec des
        variantes d'âge (frontières de 50, 58, 60, 64, 65, 70 et 75~ans),
        de composition (couples d'âges mixtes dans les deux ordres,
        enfants d'âges variés) et de frais de garde (subventionnés ou
        non) --- plus de quatre mille points de comparaison ;
  \item des \textbf{points de contrôle officiels} : les maximums publiés
        (cotisation RRQ maximale, prime RAMQ maximale, SRG maximal, etc.)
        sont vérifiés un à un contre les valeurs des organismes ;
  \item au total, près de huit cents tests automatisés, exécutés à chaque
        modification du code.
\end{itemize}

Tout écart, même d'un cent, fait échouer la suite : c'est ainsi qu'ont été
découverts --- puis modélisés fidèlement --- plusieurs comportements fins
de la référence (asymétrie du supplément du SRG, exonérations de la prime
RAMQ, contrainte temporaire de l'aide sociale, traitement du revenu de
pension avant 65~ans).

\subsection{Limites et choix de modélisation}
\label{sec:limites}

Le moteur reproduit le calculateur officiel, \emph{y compris} ses choix de
modélisation et ses idiosyncrasies, qui sont documentés plutôt que
corrigés :

\begin{itemize}
  \item \textbf{revenus simplifiés} : une seule source de revenu par
        adulte ; pas de gains en capital, de dividendes, de REER ni de
        revenus de location (section~\ref{sec:modele}) ;
  \item \textbf{frais de garde} : seuil \og jeune \fg{} à l'âge $\le 5$
        (au lieu de \og moins de 7~ans \fg{}) et plafonnement agrégé
        plutôt que par enfant (poste~6) ;
  \item \textbf{allocation-logement} : loyer imputé selon la composition,
        faute d'entrée \og loyer \fg{} ; facteur d'ajustement des
        pensions extrait de la référence (poste~10) ;
  \item \textbf{crédits de frais médicaux} (postes 12 et 18) :
        systématiquement nuls, la seule dépense médicale du modèle étant
        la prime RAMQ ; la forme anormale de la réduction du poste~12 est
        reproduite telle quelle ;
  \item \textbf{Sécurité de la vieillesse} : asymétrie du supplément
        complémentaire pour les couples mixtes, reproduite par fidélité
        (poste~17) ;
  \item \textbf{périodes} : prestations de juillet à juin assimilées à
        l'année civile ; barème RAMQ de la période tarifaire en cours au
        1\ier{}~janvier.
\end{itemize}

\subsection{Avertissement}

Cet outil est \textbf{pédagogique}. Les montants qu'il produit sont
\emph{indicatifs} : ils reposent sur des cas types simplifiés et des
données fictives, et ne tiennent pas compte de l'ensemble de la situation
réelle d'un contribuable. Ils ne constituent ni un avis fiscal, ni un avis
juridique. Pour toute décision, consulter un professionnel ou les
organismes administrateurs (Revenu Québec, Agence du revenu du Canada,
Retraite Québec, Service Canada).

\section*{Colophon}
\addcontentsline{toc}{section}{Colophon}

Ce document a été rédigé par \textbf{Claude Fable~5} (Anthropic),
assistant d'intelligence artificielle, le \textbf{10 juin 2026}, sous la
direction d'Alain Boisvert, à partir du code source du moteur de calcul et
de sa documentation de traçage. La mise en page reprend la classe
\LaTeX{} des préimpressions FAIR (Meta), telle que publiée avec
arXiv:2412.06769, adaptée au français (le logo de Meta, marque de
l'entreprise, a été retiré).

Le code source du calculateur --- moteur, application Web, suite de tests
et le présent document --- est publié à l'adresse \ghurl.
